% Generated by scripts/tikz_reviews.py using the compiled book's labels.
\pdfbookmark[0]{Chapter 9}{comparison-chapter-9}
\ReviewFigure{fig-bound-regul-variance}
{Capped inversion and Tikhonov filtering}
{figures/inverse-problems/bound-variance.png}
{figures/tikz/inverse-problems-variance-bound.pdf}
{Spectral regularization and the bound on noise amplification.}
{Interpretation to check: The left plateau is interpreted as the capped inverse 1/max(sigma,tau), rather than truncated SVD. A separate cutoff tau avoids confusing its parameter with the squared-scale Tikhonov parameter lambda. Both plots now use common axis scales and tau=sqrt(lambda): the Tikhonov maximum at sigma=sqrt(lambda) is half the capped gain. The dashed reciprocal curves have the same scale.}
{inverse-problems-variance-bound}
{Complete figure}
{9.2}
\ReviewFigure{fig-ip-bound-regul}
{The source-order dependence of the bias bound}
{figures/inverse-problems/bound-bias.png}
{figures/tikz/inverse-problems-bias-bound.pdf}
{The Tikhonov proof bounds lambda*sigma\textasciicircum{}beta/(lambda+sigma\textasciicircum{}2) uniformly over sigma>=0.}
{Corrects the handwritten maximizer: sigma*=sqrt(beta*lambda/(2-beta)), including the missing factor beta. For beta=2 the supremum lambda is approached asymptotically. For beta>2 the supremum over all sigma>=0 is infinite. The figure explicitly distinguishes this from the bounded spectrum of a fixed bounded operator. Curves show representative beta values with rescaled axes and lambda>0.}
{inverse-problems-bias-bound}
{Complete figure}
{9.3}
